\documentclass[12pt,a4paper]{article}
\usepackage[utf8]{inputenc}
\usepackage[vietnamese]{babel}
\usepackage{amsmath,amssymb,amsthm}
\usepackage{geometry}
\usepackage{enumitem}
\usepackage[most]{tcolorbox}
\usepackage{hyperref}
\usepackage{fancyhdr}
\usepackage{xcolor}
\usepackage{array}
\usepackage{booktabs}

\geometry{a4paper, margin=2.5cm, bottom=3cm}
\hypersetup{colorlinks=true, linkcolor=blue!70!black, urlcolor=blue}
\setlist[itemize]{topsep=4pt, itemsep=2pt}
\setlist[enumerate]{topsep=4pt, itemsep=2pt}

\pagestyle{fancy}
\fancyhf{}
\fancyhead[L]{\small\textit{MML -- Chương 2: Đại số tuyến tính}}
\fancyhead[R]{\small\textit{Tóm tắt tiếng Việt}}
\fancyfoot[C]{\thepage}
\renewcommand{\headrulewidth}{0.4pt}
\renewcommand{\footrulewidth}{0.4pt}

\newtcolorbox{dinhnghia}[1][]{enhanced, breakable,
  colback=blue!3!white, colframe=blue!60!black,
  fonttitle=\bfseries, title={Định nghĩa}, #1}
\newtcolorbox{dinhly}[1][]{enhanced, breakable,
  colback=green!3!white, colframe=green!50!black,
  fonttitle=\bfseries, title={Định lý / Tính chất}, #1}
\newtcolorbox{vidu}[1][]{enhanced, breakable,
  colback=violet!3!white, colframe=violet!60!black,
  fonttitle=\bfseries, title={Ví dụ}, #1}
\newtcolorbox{luuy}[1][]{enhanced, breakable,
  colback=orange!5!white, colframe=orange!60!black,
  fonttitle=\bfseries, title={Lưu ý}, #1}
\newtcolorbox{phuongphap}[1][]{enhanced, breakable,
  colback=gray!5!white, colframe=gray!50!black,
  fonttitle=\bfseries, title={Phương pháp / Thuật toán}, #1}
\newtcolorbox{tomtat}[1][]{enhanced, breakable,
  colback=teal!3!white, colframe=teal!60!black,
  fonttitle=\bfseries, title={Tóm tắt}, #1}

\title{\textbf{Chương 2: Đại số tuyến tính}\\[4pt]
\large Tóm tắt từ \textit{Mathematics for Machine Learning}\\
\normalsize Deisenroth, Faisal, Ong}
\author{}
\date{}

\begin{document}
\maketitle
\tableofcontents
\newpage

%% ============================================================
\section{Giới thiệu}

Đại số tuyến tính nghiên cứu \textbf{vector} và các quy tắc thao tác trên vector. Vector là đối tượng đặc biệt: có thể \emph{cộng với nhau} và \emph{nhân với vô hướng (scalar)} để tạo ra đối tượng cùng loại.

Các ví dụ về vector:
\begin{itemize}
  \item \textbf{Vector hình học:} Các đoạn thẳng có hướng trong không gian 2D/3D.
  \item \textbf{Đa thức:} Hai đa thức có thể cộng lại hoặc nhân với vô hướng, kết quả vẫn là đa thức.
  \item \textbf{Tín hiệu âm thanh:} Dãy số, có thể cộng và nhân vô hướng.
  \item \textbf{Phần tử của $\mathbb{R}^n$:} Bộ $n$ số thực -- đây là loại vector trọng tâm của cuốn sách.
\end{itemize}

Ý tưởng quan trọng: \textbf{tính đóng} -- tập hợp tất cả các đối tượng sinh ra từ phép cộng và nhân vô hướng từ một tập nhỏ tạo thành \emph{không gian vector}.

%% ============================================================
\section{Hệ phương trình tuyến tính (§2.1)}

Hệ phương trình tuyến tính là nền tảng của đại số tuyến tính. Dạng tổng quát:
\[
  a_{11}x_1 + a_{12}x_2 + \cdots + a_{1n}x_n = b_1
\]
\[
  \vdots
\]
\[
  a_{m1}x_1 + a_{m2}x_2 + \cdots + a_{mn}x_n = b_m
\]
trong đó $a_{ij}, b_i \in \mathbb{R}$ là hằng số đã biết, $x_j$ là các ẩn số.

Hệ phương trình này có thể viết gọn dưới dạng ma trận: $Ax = b$.

\begin{luuy}
Một hệ phương trình tuyến tính thực có thể có:
\begin{itemize}
  \item \textbf{Không có nghiệm} (vô nghiệm),
  \item \textbf{Đúng một nghiệm} (nghiệm duy nhất),
  \item \textbf{Vô số nghiệm}.
\end{itemize}
\end{luuy}

\textbf{Ý nghĩa hình học:} Trong hệ 2 ẩn, mỗi phương trình tuyến tính biểu diễn một đường thẳng. Nghiệm của hệ là giao điểm của các đường thẳng đó. Với 3 ẩn, mỗi phương trình là một mặt phẳng.

%% ============================================================
\section{Ma trận (§2.2)}

\begin{dinhnghia}
Ma trận thực $(m, n)$ là mảng $m \times n$ phần tử $a_{ij} \in \mathbb{R}$, sắp xếp theo $m$ hàng và $n$ cột:
\[
  A = \begin{pmatrix} a_{11} & \cdots & a_{1n} \\ \vdots & \ddots & \vdots \\ a_{m1} & \cdots & a_{mn} \end{pmatrix}, \quad a_{ij} \in \mathbb{R}
\]
Ký hiệu: $A \in \mathbb{R}^{m \times n}$.
\end{dinhnghia}

\subsection{Cộng và nhân ma trận}

\textbf{Cộng:} $A + B$ được tính theo từng phần tử (yêu cầu cùng kích thước).

\textbf{Nhân:} Với $A \in \mathbb{R}^{m \times n}$, $B \in \mathbb{R}^{n \times k}$, tích $C = AB \in \mathbb{R}^{m \times k}$ có:
\[
  c_{ij} = \sum_{l=1}^{n} a_{il} b_{lj}
\]

\begin{luuy}
\begin{itemize}
  \item Nhân ma trận \textbf{không giao hoán}: $AB \ne BA$ nói chung.
  \item Nhân ma trận \textbf{không phải} phép nhân theo từng phần tử (phép nhân từng phần tử gọi là tích Hadamard).
  \item Hai ma trận chỉ nhân được khi các chiều ``lân cận'' khớp nhau: $A_{m \times n} \cdot B_{n \times k} = C_{m \times k}$.
\end{itemize}
\end{luuy}

\begin{dinhly}
Tính chất phép nhân ma trận:
\begin{itemize}
  \item \textbf{Kết hợp:} $(AB)C = A(BC)$
  \item \textbf{Phân phối:} $(A+B)C = AC + BC$; $A(C+D) = AC + AD$
  \item \textbf{Ma trận đơn vị:} $I_m A = A I_n = A$
\end{itemize}
\end{dinhly}

\subsection{Ma trận nghịch đảo và chuyển vị}

\begin{dinhnghia}
\textbf{Ma trận nghịch đảo:} Với $A \in \mathbb{R}^{n \times n}$, nếu tồn tại $B$ sao cho $AB = I_n = BA$ thì $B$ gọi là nghịch đảo của $A$, ký hiệu $A^{-1}$. $A$ được gọi là khả nghịch (regular/invertible/nonsingular). Nếu không tồn tại, $A$ là suy biến (singular).
\end{dinhnghia}

\begin{dinhnghia}
\textbf{Ma trận chuyển vị:} Với $A \in \mathbb{R}^{m \times n}$, ma trận $B \in \mathbb{R}^{n \times m}$ với $b_{ij} = a_{ji}$ gọi là chuyển vị của $A$, ký hiệu $A^\top$.
\end{dinhnghia}

\begin{dinhly}
Các tính chất quan trọng:
\begin{itemize}
  \item $AA^{-1} = I = A^{-1}A$
  \item $(AB)^{-1} = B^{-1}A^{-1}$
  \item $(A^\top)^\top = A$
  \item $(AB)^\top = B^\top A^\top$
  \item $(A + B)^\top = A^\top + B^\top$
\end{itemize}
\end{dinhly}

\begin{dinhnghia}
\textbf{Ma trận đối xứng:} $A \in \mathbb{R}^{n \times n}$ là đối xứng nếu $A = A^\top$.
\end{dinhnghia}

%% ============================================================
\section{Giải hệ phương trình tuyến tính (§2.3)}

\subsection{Nghiệm riêng và nghiệm tổng quát}

Để giải $Ax = b$, ta tìm:
\begin{enumerate}
  \item \textbf{Nghiệm riêng (particular solution):} Một nghiệm cụ thể thỏa $Ax = b$.
  \item \textbf{Nghiệm thuần nhất (homogeneous solution):} Toàn bộ nghiệm của $Ax = 0$.
  \item \textbf{Nghiệm tổng quát:} Tổng của nghiệm riêng và tất cả nghiệm thuần nhất.
\end{enumerate}

\subsection{Các phép biến đổi sơ cấp}

Ba phép biến đổi giữ nguyên tập nghiệm:
\begin{itemize}
  \item Hoán vị hai phương trình (hai hàng).
  \item Nhân một phương trình (hàng) với hằng số $\lambda \ne 0$.
  \item Cộng hai phương trình (hàng) vào nhau.
\end{itemize}

\subsection{Dạng bậc thang hàng (Row-Echelon Form -- REF)}

\begin{dinhnghia}
Ma trận ở dạng \textbf{bậc thang hàng (REF)} khi:
\begin{itemize}
  \item Tất cả hàng gồm toàn số 0 nằm ở dưới cùng.
  \item Phần tử khác 0 đầu tiên từ trái của mỗi hàng (gọi là \emph{phần tử chốt -- pivot}) nằm nghiêm ngặt bên phải pivot của hàng trên nó.
\end{itemize}
\end{dinhnghia}

\begin{dinhnghia}
\textbf{Dạng bậc thang hàng rút gọn (Reduced REF -- RREF):} REF thỏa thêm: mỗi pivot bằng 1 và là phần tử khác 0 duy nhất trong cột của nó.
\end{dinhnghia}

\begin{phuongphap}[title={Khử Gauss (Gaussian Elimination)}]
Thuật toán dùng các phép biến đổi sơ cấp để đưa hệ về RREF, từ đó đọc ra nghiệm:
\begin{enumerate}
  \item Đưa ma trận về REF.
  \item Tiếp tục đưa về RREF.
  \item Các biến ứng với cột pivot gọi là \emph{biến cơ bản}; các biến còn lại là \emph{biến tự do}.
  \item Đặt biến tự do là tham số $\lambda_i \in \mathbb{R}$, biểu diễn nghiệm tổng quát.
\end{enumerate}
\end{phuongphap}

\subsection{Tính ma trận nghịch đảo bằng khử Gauss}

Để tính $A^{-1}$ với $A \in \mathbb{R}^{n \times n}$, ta đưa ma trận mở rộng $[A \mid I_n]$ về RREF:
\[
  [A \mid I_n] \;\xrightarrow{\text{Gauss}}\; [I_n \mid A^{-1}]
\]

%% ============================================================
\section{Không gian vector (§2.4)}

\subsection{Nhóm (Group)}

\begin{dinhnghia}
Cho tập $G$ và phép toán $\otimes : G \times G \to G$. $(G, \otimes)$ là \textbf{nhóm} nếu thỏa:
\begin{enumerate}
  \item \textbf{Đóng:} $\forall x, y \in G: x \otimes y \in G$
  \item \textbf{Kết hợp:} $(x \otimes y) \otimes z = x \otimes (y \otimes z)$
  \item \textbf{Phần tử đơn vị:} $\exists e \in G: x \otimes e = x$ và $e \otimes x = x$
  \item \textbf{Phần tử nghịch đảo:} $\forall x \in G, \exists y \in G: x \otimes y = e$ và $y \otimes x = e$
\end{enumerate}
Nếu thêm $x \otimes y = y \otimes x$ với mọi $x, y$, thì $(G, \otimes)$ là \textbf{nhóm Abel (giao hoán)}.
\end{dinhnghia}

\begin{vidu}
$(\mathbb{Z}, +)$ là nhóm Abel. $(\mathbb{R} \setminus \{0\}, \cdot)$ là nhóm Abel. $(\mathbb{R}^{n \times n}, \cdot)$ với các ma trận khả nghịch (nhóm tuyến tính tổng quát $GL(n, \mathbb{R})$) là nhóm (không giao hoán).
\end{vidu}

\subsection{Không gian vector}

\begin{dinhnghia}
\textbf{Không gian vector thực} là bộ $V = (V, +, \cdot)$ gồm tập $V$, phép cộng $+ : V \times V \to V$ và phép nhân vô hướng $\cdot : \mathbb{R} \times V \to V$, thỏa:
\begin{itemize}
  \item $(V, +)$ là nhóm Abel.
  \item Phân phối: $\lambda(x + y) = \lambda x + \lambda y$; $(\lambda + \psi)x = \lambda x + \psi x$
  \item Kết hợp (với vô hướng): $\lambda(\psi x) = (\lambda \psi)x$
  \item Đơn vị: $1 \cdot x = x$
\end{itemize}
Các phần tử của $V$ gọi là \textbf{vector}; $\mathbb{R}$ là trường vô hướng.
\end{dinhnghia}

\begin{dinhnghia}
\textbf{Không gian con (Subspace):} Tập $U \subseteq V$ không rỗng gọi là không gian con của $V$ nếu $U$ đóng với phép cộng và nhân vô hướng, tức $U$ tự thành một không gian vector.
\end{dinhnghia}

%% ============================================================
\section{Độc lập tuyến tính (§2.5)}

\begin{dinhnghia}
Tập hợp vector $\{x_1, \ldots, x_k\} \subset V$ gọi là \textbf{độc lập tuyến tính} nếu phương trình:
\[
  \lambda_1 x_1 + \lambda_2 x_2 + \cdots + \lambda_k x_k = 0
\]
chỉ có nghiệm duy nhất $\lambda_1 = \lambda_2 = \cdots = \lambda_k = 0$.
Ngược lại, tập hợp gọi là \textbf{phụ thuộc tuyến tính}.
\end{dinhnghia}

\begin{dinhly}
Tính chất:
\begin{itemize}
  \item Nếu một trong các vector là vector không, tập hợp phụ thuộc tuyến tính.
  \item Nếu một vector là tổ hợp tuyến tính của các vector còn lại, tập hợp phụ thuộc tuyến tính.
  \item $m$ tổ hợp tuyến tính của $k$ vector sẽ phụ thuộc tuyến tính nếu $m > k$.
\end{itemize}
\end{dinhly}

\begin{phuongphap}[title={Kiểm tra độc lập tuyến tính}]
Đặt các vector làm cột của ma trận $A$, sau đó dùng khử Gauss đưa về REF. Nếu mọi cột đều là cột pivot $\Rightarrow$ độc lập tuyến tính. Nếu có cột không phải pivot $\Rightarrow$ phụ thuộc tuyến tính.
\end{phuongphap}

%% ============================================================
\section{Cơ sở và Hạng (§2.6)}

\subsection{Tập sinh và cơ sở}

\begin{dinhnghia}
Cho không gian vector $V$ và $A = \{x_1, \ldots, x_k\} \subseteq V$. Nếu mọi $v \in V$ đều biểu diễn được dưới dạng tổ hợp tuyến tính của các phần tử trong $A$, thì $A$ gọi là \textbf{tập sinh} của $V$. Bao tuyến tính (span) của $A$:
\[
  \mathrm{span}[A] = \left\{\sum_{i=1}^k \lambda_i x_i \;\Big|\; \lambda_i \in \mathbb{R}\right\}
\]
\end{dinhnghia}

\begin{dinhnghia}
\textbf{Cơ sở (Basis):} Tập sinh $A$ của $V$ gọi là \textbf{cơ sở} nếu $A$ là tập sinh tối tiểu (không có tập sinh nào nhỏ hơn), hay tương đương, $A$ là tập độc lập tuyến tính lớn nhất trong $V$.
\end{dinhnghia}

\begin{dinhly}
Các mệnh đề tương đương:
\begin{itemize}
  \item $B$ là cơ sở của $V$.
  \item $B$ là tập sinh tối tiểu.
  \item $B$ là tập độc lập tuyến tính lớn nhất.
  \item Mọi $x \in V$ đều biểu diễn được duy nhất qua các phần tử của $B$.
\end{itemize}
\end{dinhly}

\begin{luuy}
Cơ sở tắc (canonical basis) của $\mathbb{R}^n$: $\{e_1, e_2, \ldots, e_n\}$, trong đó $e_i$ là vector có 1 ở vị trí thứ $i$ và 0 ở các vị trí còn lại. Mỗi không gian vector có nhiều cơ sở khác nhau, nhưng tất cả có cùng số phần tử.
\end{luuy}

\begin{dinhnghia}
\textbf{Số chiều (Dimension):} Số chiều $\dim(V)$ của không gian vector $V$ là số lượng vector cơ sở. Nếu $U \subseteq V$ là không gian con thì $\dim(U) \le \dim(V)$.
\end{dinhnghia}

\subsection{Hạng (Rank)}

\begin{dinhnghia}
\textbf{Hạng} của ma trận $A \in \mathbb{R}^{m \times n}$, ký hiệu $\mathrm{rk}(A)$, là số cột (hoặc hàng) độc lập tuyến tính tối đa của $A$.
\end{dinhnghia}

\begin{dinhly}
Tính chất quan trọng của hạng:
\begin{itemize}
  \item $\mathrm{rk}(A) = \mathrm{rk}(A^\top)$ (hạng cột = hạng hàng).
  \item $A \in \mathbb{R}^{n \times n}$ khả nghịch $\Leftrightarrow$ $\mathrm{rk}(A) = n$.
  \item Hệ $Ax = b$ có nghiệm $\Leftrightarrow$ $\mathrm{rk}(A) = \mathrm{rk}([A \mid b])$.
  \item Không gian nghiệm của $Ax = 0$ có số chiều $n - \mathrm{rk}(A)$ (gọi là không gian nhân -- kernel/null space).
  \item Ma trận \emph{full rank} khi $\mathrm{rk}(A) = \min(m, n)$.
\end{itemize}
\end{dinhly}

%% ============================================================
\section{Ánh xạ tuyến tính (§2.7)}

\begin{dinhnghia}
Cho hai không gian vector $V, W$. Ánh xạ $\Phi : V \to W$ gọi là \textbf{ánh xạ tuyến tính} (hay đồng cấu không gian vector) nếu:
\[
  \forall x, y \in V,\; \forall \lambda, \psi \in \mathbb{R}:\quad \Phi(\lambda x + \psi y) = \lambda\Phi(x) + \psi\Phi(y)
\]
\end{dinhnghia}

Các trường hợp đặc biệt:
\begin{itemize}
  \item \textbf{Đẳng cấu (Isomorphism):} $\Phi : V \to W$ tuyến tính và song ánh (bijective).
  \item \textbf{Nội cấu (Endomorphism):} $\Phi : V \to V$ tuyến tính.
  \item \textbf{Tự đẳng cấu (Automorphism):} $\Phi : V \to V$ tuyến tính và song ánh.
\end{itemize}

\begin{dinhly}[title={Định lý về đẳng cấu}]
Hai không gian vector hữu hạn chiều $V$ và $W$ đẳng cấu nhau khi và chỉ khi $\dim(V) = \dim(W)$. Điều này có nghĩa là hai không gian vector cùng chiều về bản chất là ``như nhau'' -- chúng có thể biến đổi qua lại mà không mất thông tin.
\end{dinhly}

\subsection{Biểu diễn ma trận của ánh xạ tuyến tính}

\begin{dinhnghia}
Cho không gian vector $V, W$ với cơ sở có thứ tự $B = (b_1, \ldots, b_n)$ và $C = (c_1, \ldots, c_m)$. Ánh xạ tuyến tính $\Phi : V \to W$ tương ứng với \textbf{ma trận chuyển đổi} $A_\Phi \in \mathbb{R}^{m \times n}$, trong đó cột $j$ của $A_\Phi$ chứa tọa độ của $\Phi(b_j)$ theo $C$.
\end{dinhnghia}

Nếu $\hat{x}$ là vector tọa độ của $x \in V$ theo $B$ và $\hat{y}$ là vector tọa độ của $\Phi(x) \in W$ theo $C$, thì:
\[
  \hat{y} = A_\Phi \hat{x}
\]

\begin{luuy}
Khi làm việc với ma trận, cần phân biệt hai vai trò:
\begin{itemize}
  \item Ma trận biểu diễn một \textbf{ánh xạ tuyến tính}.
  \item Ma trận là \textbf{tập hợp các vector} (mỗi cột là một vector).
\end{itemize}
\end{luuy}

\subsection{Thay đổi cơ sở (Basis Change)}

Khi thay đổi cơ sở trong $V$ và $W$, ma trận chuyển đổi thay đổi theo:
\[
  \tilde{A}_\Phi = T^{-1} A_\Phi S
\]
trong đó $S$ là ma trận chuyển đổi tọa độ từ cơ sở mới sang cơ sở cũ trong $V$, và $T$ là ma trận tương ứng trong $W$.

\begin{dinhnghia}
\textbf{Tương đương:} $A, \tilde{A} \in \mathbb{R}^{m \times n}$ tương đương nếu $\exists$ ma trận khả nghịch $S \in \mathbb{R}^{n \times n}$, $T \in \mathbb{R}^{m \times m}$: $\tilde{A} = T^{-1}AS$.

\textbf{Tương tự (Similar):} $A, \tilde{A} \in \mathbb{R}^{n \times n}$ tương tự nếu $\exists$ ma trận khả nghịch $S \in \mathbb{R}^{n \times n}$: $\tilde{A} = S^{-1}AS$.

Ma trận tương tự thì luôn tương đương; điều ngược lại không nhất thiết đúng.
\end{dinhnghia}

\subsection{Ảnh và Nhân (Image \& Kernel)}

\begin{dinhnghia}
Cho ánh xạ tuyến tính $\Phi : V \to W$:
\begin{itemize}
  \item \textbf{Nhân (Kernel / Null Space):}
  \[
    \ker(\Phi) = \{v \in V : \Phi(v) = 0_W\}
  \]
  Tập hợp các vector trong $V$ bị ánh xạ về $0_W$.
  \item \textbf{Ảnh (Image / Range):}
  \[
    \mathrm{Im}(\Phi) = \{\Phi(v) : v \in V\}
  \]
  Tập hợp tất cả các vector trong $W$ có thể đạt được từ $\Phi$.
\end{itemize}
\end{dinhnghia}

\begin{dinhly}
Tính chất:
\begin{itemize}
  \item $\ker(\Phi)$ luôn chứa $0_V$ (không bao giờ rỗng).
  \item $\mathrm{Im}(\Phi) \subseteq W$ là không gian con của $W$.
  \item $\ker(\Phi) \subseteq V$ là không gian con của $V$.
  \item $\Phi$ đơn ánh (injective) $\Leftrightarrow$ $\ker(\Phi) = \{0\}$.
\end{itemize}
\end{dinhly}

\begin{dinhly}[title={Định lý hạng -- nhân (Rank-Nullity Theorem)}]
Với ánh xạ tuyến tính $\Phi : V \to W$:
\[
  \dim(\ker(\Phi)) + \dim(\mathrm{Im}(\Phi)) = \dim(V)
\]
\end{dinhly}

\begin{luuy}
Với $A \in \mathbb{R}^{m \times n}$ và ánh xạ $\Phi : \mathbb{R}^n \to \mathbb{R}^m$, $x \mapsto Ax$:
\begin{itemize}
  \item $\mathrm{Im}(\Phi) = \mathrm{span}[a_1, \ldots, a_n]$ (bao tuyến tính của các cột của $A$) -- còn gọi là \emph{không gian cột}.
  \item $\ker(\Phi)$ là nghiệm tổng quát của $Ax = 0$.
  \item $\mathrm{rk}(A) = \dim(\mathrm{Im}(\Phi))$.
\end{itemize}
\end{luuy}

%% ============================================================
\section{Không gian Affine (§2.8)}

\begin{dinhnghia}
Cho không gian vector $V$, $x_0 \in V$ và không gian con $U \subseteq V$. Tập hợp:
\[
  L = x_0 + U := \{x_0 + u : u \in U\}
\]
gọi là \textbf{không gian con affine} của $V$. $U$ gọi là \emph{không gian hướng}, $x_0$ gọi là \emph{điểm neo (support point)}.
\end{dinhnghia}

\begin{luuy}
Không gian affine \textbf{không} phải không gian con (vector subspace) nếu $x_0 \notin U$, vì không chứa $0$.
\end{luuy}

Ví dụ về không gian affine trong $\mathbb{R}^n$:
\begin{itemize}
  \item \textbf{Đường thẳng (không gian 1 chiều):} $y = x_0 + \lambda b_1$, $\lambda \in \mathbb{R}$.
  \item \textbf{Mặt phẳng (không gian 2 chiều):} $y = x_0 + \lambda_1 b_1 + \lambda_2 b_2$.
  \item \textbf{Siêu phẳng (hyperplane, không gian $(n-1)$ chiều):} $y = x_0 + \sum_{i=1}^{n-1} \lambda_i b_i$.
\end{itemize}

\begin{dinhnghia}
\textbf{Ánh xạ affine:} Với ánh xạ tuyến tính $\Phi : V \to W$ và $a \in W$, ánh xạ:
\[
  \varphi : V \to W,\quad x \mapsto a + \Phi(x)
\]
gọi là ánh xạ affine. Vector $a$ gọi là \emph{vector dịch chuyển (translation vector)}.
\end{dinhnghia}

\begin{luuy}
Mọi ánh xạ affine đều là hợp thành của một ánh xạ tuyến tính và một phép tịnh tiến. Hợp thành của hai ánh xạ affine cũng là affine.
\end{luuy}

%% ============================================================
\section{Tóm tắt chương}

\begin{tomtat}
\begin{itemize}
  \item \textbf{Vector} là đối tượng có thể cộng và nhân vô hướng, tạo ra đối tượng cùng loại.
  \item \textbf{Hệ phương trình tuyến tính} $Ax = b$ có thể giải bằng khử Gauss đưa về dạng RREF.
  \item \textbf{Ma trận} biểu diễn ánh xạ tuyến tính hoặc tập hợp vector; có thể cộng, nhân, lấy nghịch đảo và chuyển vị.
  \item \textbf{Không gian vector} là nơi vector ``sống''; được đặc trưng bởi các tiên đề về phép cộng và nhân vô hướng.
  \item \textbf{Độc lập tuyến tính}: tập vector không thể biểu diễn lẫn nhau qua tổ hợp tuyến tính.
  \item \textbf{Cơ sở}: tập sinh tối tiểu; số phần tử = số chiều không gian.
  \item \textbf{Hạng} của ma trận = số hàng/cột độc lập tuyến tính tối đa.
  \item \textbf{Ánh xạ tuyến tính} bảo toàn cấu trúc không gian vector; biểu diễn bằng ma trận.
  \item \textbf{Nhân/Ảnh} của ánh xạ tuyến tính liên hệ với nhau qua định lý hạng-nhân.
  \item \textbf{Không gian affine}: không gian vector bị dịch chuyển (không qua gốc tọa độ).
\end{itemize}
\end{tomtat}

\end{document}
